% Source: https://artofproblemsolving.com/community/c3h1374695_mathcounts_problem
% Stable identity: aops-topic-1374695-q01
If $n=p_1^{e_1}\cdots p_k^{e_k}$, then $n$ has $(e_1+1)\cdots(e_k+1)$ positive factors. The ways to write $12$ as such a product give the possible shapes ($p,q,r$ distinct primes):
\begin{itemize}
\item $p^{11}$: at least $2^{11}=2048$, too large.
\item $p^5q$: $2^5\cdot3=96$ works; $2^5\cdot5=160$ and $3^5\cdot2=486$ are too large.
\item $p^3q^2$: $2^3\cdot3^2=72$ works; $2^3\cdot5^2=200$, $3^3\cdot2^2=108$ are too large.
\item $p^2qr$: $2^2\cdot3\cdot5=60$, $2^2\cdot3\cdot7=84$, $3^2\cdot2\cdot5=90$ work; the next ones, $2^2\cdot3\cdot11=132$, $2^2\cdot5\cdot7=140$, $3^2\cdot2\cdot7=126$, $5^2\cdot2\cdot3=150$, are too large.
\end{itemize}
The two-digit numbers are $60,72,84,90,96$, with sum $\boxed{402}$.
